\subsection{Heisenberg Representation of Quantum Circuits}

Up to now, we have described quantum computation using the \textbf{Schrödinger representation}. In this approach, the \hl{quantum state evolves over time}, while the observables remain fixed.

\highspace
However, there is an alternative but completely equivalent description called the \textbf{Heisenberg representation}. In this approach, the \hl{quantum state remains fixed}, while the observables evolve instead.

\highspace
Both descriptions produce exactly the same physical predictions and measurement outcomes. They are just different mathematical viewpoints of the same underlying physics.

\highspace
\begin{flushleft}
    \textcolor{Green3}{\faIcon{book} \textbf{Schrödinger Representation}}
\end{flushleft}
In the \definition{Schrödinger Representation}, a quantum gate $U$ acts directly on the state:
\begin{equation}
    \ket{\psi_1}
    =
    U\ket{\psi_0}
\end{equation}
Where:
\begin{itemize}
    \item $\ket{\psi_0}$ is the initial state;
    \item $U$ is the unitary operator representing the quantum circuit;
    \item $\ket{\psi_1}$ is the final state.
\end{itemize}
The \hl{observable does not change.} For example, if we measure in the computational basis, we use the Pauli-Z observable $Z$, which \textbf{remains the same before and after applying the circuit}. The expectation value after applying the circuit is:
\begin{equation*}
    \bra{\psi_1} Z \ket{\psi_1}
\end{equation*}
In this case:
\begin{itemize}
    \item The \textbf{state changes};
    \item The \textbf{observable stays fixed}.
\end{itemize}

\highspace
\begin{flushleft}
    \textcolor{Green3}{\faIcon{book} \textbf{Heisenberg Representation}}
\end{flushleft}
In the \definition{Heisenberg Representation}, the state remains unchanged $\ket{\psi_0}$.
Instead, the \textbf{observable is transformed} according to:
\begin{equation}\label{eq:heisenberg-evolution}
    Z
    \;\longrightarrow\;
    U^\dagger Z U
\end{equation}
Where:
\begin{itemize}
    \item $U^{\dagger}$ is the conjugate transpose of the gate $U$;
    \item $Z$ is the original observable;
    \item $U$ is the gate itself.
\end{itemize}
The expectation value is then computed as:
\begin{equation*}
    \bra{\psi_0}
    U^\dagger Z U
    \ket{\psi_0}
\end{equation*}
In this case:
\begin{itemize}
    \item The \textbf{state stays fixed};
    \item The \textbf{observable changes}.
\end{itemize}

\highspace
\begin{flushleft}
    \textcolor{Green3}{\faIcon{question-circle} \textbf{Why are the two representations equivalent?}}
\end{flushleft}
Starting from the Schrödinger representation:
\begin{equation*}
    \ket{\psi_1}
    =
    U\ket{\psi_0}
\end{equation*}
The corresponding bra is:
\begin{equation*}
    \bra{\psi_1}
    =
    \bra{\psi_0} U^\dagger
\end{equation*}
Substituting into the expectation value:
\begin{align*}
    \bra{\psi_1} Z \ket{\psi_1}
    &=
    \left(\bra{\psi_0} U^\dagger\right)
    Z
    \left(U\ket{\psi_0}\right)
    \\
    &=
    \bra{\psi_0}
    U^\dagger Z U
    \ket{\psi_0}
\end{align*}
Therefore (and this is the key point), the \hl{expectation value computed in the Schrödinger representation is exactly equal to the expectation value computed in the Heisenberg representation} (see \autoref{eq:heisenberg-evolution}):
\begin{equation}\label{eq:schrodinger-heisenberg-equivalence}
    \bra{\psi_1} Z \ket{\psi_1}
    =
    \bra{\psi_0}
    U^\dagger Z U
    \ket{\psi_0}
\end{equation}
This proves that both representations yield exactly the same measurement result.

\highspace
\begin{flushleft}
    \textcolor{Green3}{\faIcon{balance-scale} \textbf{Physical Interpretation}}
\end{flushleft}
The Schrödinger and Heisenberg representations are two different mathematical viewpoints of the same physical process.

\highspace
In the \textbf{Schrödinger} representation:
\begin{itemize}
    \item We imagine the quantum state moving through the circuit;
    \item Gates transform the state;
    \item Measurements are performed using fixed observables.
\end{itemize}
In the \textbf{Heisenberg} representation:
\begin{itemize}
    \item We keep the quantum state fixed;
    \item Gates transform the observables;
    \item Measurements are performed using transformed observables.
\end{itemize}

\highspace
\begin{flushleft}
    \textcolor{Green3}{\faIcon{check-circle} \textbf{Why is the Heisenberg representation useful?}}
\end{flushleft}
The Heisenberg representation is useful because it lets us answer a very important question: \emph{\textbf{what observable is this quantum circuit actually measuring?}}

\highspace
For example, in the Schrödinger representation, suppose we have an initial state $\ket{\psi}$, a quantum circuit $H$ (the Hadamard gate, \autopageref{subsubsec:hadamard-gate-h}), and we measure in the computational basis using the observable $Z$ (\autopageref{sec:pauli-z-phase-flip-gate}). \emph{But what does this circuit actually measure?}

\highspace
To find out, we can switch to the Heisenberg representation and compute the transformed observable:
\begin{equation*}
    H^\dagger Z H
\end{equation*}
This will give us the new observable that corresponds to the measurement performed by the circuit $H$. In this case, we find that:
\begin{equation*}
    H^\dagger Z H =
    \dfrac{1}{2} \left(
        \begin{bmatrix}
            1 & 1 \\
            1 & -1
        \end{bmatrix}
        \begin{bmatrix}
            1 & 0 \\
            0 & -1
        \end{bmatrix}
        \begin{bmatrix}
            1 & 1 \\
            1 & -1
        \end{bmatrix}
    \right) =
    \dfrac{1}{2}
    \begin{bmatrix}
        0 & 2 \\
        2 & 0
    \end{bmatrix} =
    \begin{bmatrix}
        0 & 1 \\
        1 & 0
    \end{bmatrix} = X
\end{equation*}
This means that the circuit $H$ actually measures the observable $X$ (the Pauli-X gate, \autopageref{sec:pauli-x-not-gate}), not $Z$. In other words, applying the Hadamard gate before measuring in the computational basis effectively changes our measurement from $Z$ to $X$.

\highspace
So, without the Heisenberg representation, we only know the sequence of gates. With the Heisenberg representation, we immediately understand: ``\emph{this circuit is actually measuring $X$!}'' This is the real physical meaning of the circuit, and it is crucial for understanding how quantum algorithms work.

\begin{examplebox}[: Analogy with Arithmetic Expressions]
    Suppose someone gives us:
    \begin{equation*}
        \left(3 + 5\right) \times 2 - 16
    \end{equation*}
    We could compute it step by step. Or we could simplify the whole expression to: $0$. The Heisenberg representation does something similar:
    \begin{itemize}
        \item Instead of following every state transformation,
        \item It simplifies the entire circuit into an equivalent measurement.
    \end{itemize}
\end{examplebox}

\noindent
Suppose a circuit contains many gates. In the \hl{Schrödinger picture, we must track the state after each gate.} In the \hl{Heisenberg picture, we ask only: \emph{what observable is measured at the end?}} Sometimes the answer is a single operator such as $X$ or $Z$. Sometimes it is a more complicated operator that involves multiple qubits. But in either case, we get a clear understanding of what the circuit does without having to track the state at every step.